%% Section 3: The Inert Clipping Problem
%% To be included in main paper via %% Section 3: The Inert Clipping Problem
%% To be included in main paper via %% Section 3: The Inert Clipping Problem
%% To be included in main paper via %% Section 3: The Inert Clipping Problem
%% To be included in main paper via \input{section3}

\section{The Inert Clipping Problem}
\label{sec:inert-clipping}

PPO's clipped surrogate objective is the central trust-region mechanism inherited by GRPO.
In this section, we demonstrate that this mechanism is empirically inert across all standard training configurations we tested, identify two independent root causes, and show how this diagnosis unifies a series of recent algorithmic redesigns.

\subsection{Measuring Clipping Activity}
\label{sec:measuring-clipping}

We instrument the PPO loss to track \texttt{pg\_clipfrac}, the fraction of tokens at which the clipped and unclipped objectives differ---i.e., the fraction of tokens where clipping actively modifies the gradient.
A value of $\texttt{pg\_clipfrac} = 0$ means the clip operator in Equation~\ref{eq:ppo-clip} is never binding: for every token, $\text{clip}(r_t, 1-\varepsilon, 1+\varepsilon) \cdot \hat{A}_t = r_t \cdot \hat{A}_t$, and the trust-region constraint contributes nothing to the optimization.

\paragraph{Setup.}
We train Qwen2.5-Math-1.5B on the DAPO-Math-17K dataset under standard GRPO configurations, resuming from a checkpoint at iteration~799.
All runs use the default clip threshold $\varepsilon = 0.2$.
We additionally track the Effective Sample Size (ESS), defined as
\begin{equation}
\label{eq:ess}
\text{ESS} = \frac{1}{\sum_{i} w_i^2}, \quad w_i = \frac{r_i}{\sum_j r_j},
\end{equation}
where $r_i = \pi_\theta(a_i \mid s) / \pi_{\text{old}}(a_i \mid s)$ are the per-token importance ratios.
ESS measures how many effective independent samples remain after importance weighting; $\text{ESS} \approx 1.0$ indicates near-perfect on-policy training.

\paragraph{PPO epochs do not activate clipping.}
At $\text{lr} = 10^{-6}$, we vary the number of PPO epochs $K \in \{1, 3, 5\}$.
In standard GRPO, multiple PPO epochs recompute log-probabilities at the start of each epoch, resetting the importance ratio denominator.
Across all three settings, $\texttt{pg\_clipfrac} = 0.0$ exactly, and $\text{ESS} = 0.999$.
Increasing the number of optimization passes over the same rollout data does not induce any clipping activity.

\paragraph{Data reuse does not activate clipping.}
We increase the learning rate to $10^{-5}$ and set \texttt{--num-steps-per-rollout}$\,{=}\,8$, which partitions each rollout into 8 non-overlapping training chunks processed sequentially \emph{without} recomputing log-probabilities between steps.
This is the most aggressive data-reuse configuration in our experimental grid.
Even here, $\texttt{pg\_clipfrac} = 0.0$ at every step position, and ESS remains above $0.999$.

\paragraph{Comparison with prior observations.}
Our finding of exactly zero clipping is consistent with, but more extreme than, prior reports.
\citet{chen2025iclr} measure $\texttt{pg\_clipfrac} < 0.2\%$ on Qwen2.5-Math-7B at $\text{lr} = 5 \times 10^{-7}$, attributing the low rate to conservative hyperparameters.
\citet{prosperity2025} report $\texttt{pg\_clipfrac} = 0.05\%$ at staleness $s = 0$ in their multi-policy framework.
Neither work investigates why the rate is near zero or identifies the architectural cause.
The small residual clipping they observe likely reflects numerical differences in their pipelines (larger models, different checkpoint states, or floating-point accumulation effects), but the qualitative conclusion is the same: \emph{clipping is essentially inactive in standard GRPO training}.


\subsection{Root Cause: Two Independent Mechanisms}
\label{sec:root-cause}

We identify two independent mechanisms that jointly render PPO clipping inert.
Their independence is important: removing either one alone is sufficient to produce nonzero clipping activity.

\subsubsection{Cause 1: Architectural Disconnection}
\label{sec:cause-architectural}

The PPO clipped surrogate objective computes the importance ratio as
\begin{equation}
\label{eq:importance-ratio}
r_t(\theta) = \frac{\pi_\theta(a_t \mid s_t)}{\pi_{\text{old}}(a_t \mid s_t)},
\end{equation}
where, in classical PPO, $\pi_{\text{old}}$ represents the \emph{behavior policy} that generated the data.
In GRPO as implemented in standard training frameworks, the denominator uses \texttt{batch["log\_probs"]}---log-probabilities \emph{recomputed} from the current model at the start of each training step---rather than \texttt{batch["rollout\_log\_probs"]}, the log-probabilities recorded at rollout generation time.

This architectural choice means that both numerator and denominator of $r_t(\theta)$ track the same policy (or very nearly so, differing only by the gradient updates within a single step).
Consequently, $r_t(\theta) \approx 1$ regardless of how far the current policy $\pi_\theta$ has drifted from the policy $\pi_{\text{rollout}}$ that actually generated the training data.
The importance ratio is \emph{disconnected} from data staleness: the PPO loss cannot detect, and therefore cannot constrain, policy drift relative to the data-generating distribution.

\subsubsection{Cause 2: Conservative Learning Rates}
\label{sec:cause-lr}

Standard GRPO training uses learning rates in the range $10^{-6}$ to $10^{-5}$.
At these rates, the per-step change in log-probabilities is small enough that even a correctly computed importance ratio (using stale rollout log-probs) would rarely exceed the clip threshold $1 + \varepsilon = 1.2$.
This provides a second, independent barrier to clipping activation: the policy simply does not move fast enough within a single rollout's training steps for the trust region to bind.

\subsubsection{Evidence for Independence}
\label{sec:evidence-independence}

We isolate the two causes by introducing a \texttt{--use-rollout-logprobs} flag that replaces the denominator of $r_t(\theta)$ with the stale rollout log-probabilities, removing Cause~1 while leaving Cause~2 intact.
At $\text{lr} = 10^{-5}$ with \texttt{--num-steps-per-rollout}$\,{=}\,8$:

\begin{itemize}
    \item \textbf{With stale log-probs} (Cause~1 removed): We observe a clear ESS gradient within each rollout, declining from $0.999$ at step~1 to $0.988$ at step~8.
    The variable \texttt{pg\_clipfrac} is nonzero at \emph{all} step positions, rising from $0.001$ (step~1) to $0.024$ (step~8) as accumulated policy drift pushes more tokens past the clip threshold.
    Over 218 rollouts, \texttt{pg\_clipfrac} is consistently positive.

    \item \textbf{Without stale log-probs} (Cause~1 present): Under identical hyperparameters, $\texttt{pg\_clipfrac} = 0.0$ at every step position regardless of $N$.
    The recomputed denominator masks all policy drift.
\end{itemize}

This contrast establishes that Cause~1 and Cause~2 are independent.
Removing the architectural disconnection (Cause~1) is \emph{sufficient} to activate clipping even at the conservative learning rate of $10^{-5}$---the rate merely modulates the \emph{magnitude} of clipping activity, not its presence.
Conversely, increasing the learning rate alone (without fixing the log-prob reference) cannot activate clipping, because the recomputed denominator absorbs all policy changes.


\subsection{Implications for Recent Algorithmic Redesigns}
\label{sec:implications}

Our diagnosis provides a unifying explanation for a striking pattern in recent work: at least six independent papers have redesigned or removed GRPO's clipping mechanism, each motivated by different empirical observations, yet all responding to the same underlying pathology.
Table~\ref{tab:related-redesigns} summarizes how each algorithm's design choice relates to our two-cause diagnosis.

\begin{table*}[t]
\centering
\caption{Recent clipping redesigns interpreted through the lens of inert clipping.
Each modification addresses a symptom of the same root cause: the PPO importance ratio is architecturally near unity, rendering clipping inactive.
``Fresh log-probs'' indicates the algorithm retains the recomputed denominator (Cause~1 unresolved).}
\label{tab:related-redesigns}
\small
\begin{tabular}{@{}p{2.0cm}p{4.5cm}p{7.5cm}@{}}
\toprule
\textbf{Algorithm} & \textbf{Modification} & \textbf{Our Explanation} \\
\midrule
DAPO \newline \citep{dapo2025}
& Asymmetric clip bounds ($\varepsilon_{\text{low}} = 0.2$, $\varepsilon_{\text{high}} = 0.28$)
& Presupposes active clipping.
  If $\texttt{pg\_clipfrac} = 0$, Clip-Higher has no tokens to asymmetrically clip.
  Their reported gains (+8 AIME) are attributable to Dynamic Sampling, not the clip modification (+2 AIME). \\
\addlinespace
GSPO \newline \citep{gspo2025}
& Sequence-level importance sampling replacing token-level
& Retains fresh log-probs; Cause~1 persists.
  Their finding that GSPO clips $100\times$ more tokens than standard GRPO directly supports our diagnosis: GRPO clips almost nothing. \\
\addlinespace
CFPO \newline \citep{cfpo2025}
& Quadratic TV-divergence penalty replacing clip
& Their penalty term is also approximately zero when $r_t \approx 1$.
  The zero-gradient problem they address only manifests after ratios deviate from unity---which requires resolving our upstream Cause~1 first. \\
\addlinespace
ASPO \newline \citep{aspo2025}
& Inverted IS ratios for positive-advantage tokens
& They report average IS weight $\approx 1.0004$ but do not investigate this observation.
  Their gradient-direction fix has negligible magnitude when $r_t \approx 1$; the effective update is dominated by the advantage term alone. \\
\addlinespace
RGRA \newline \citep{rgra2025}
& Remove clipping entirely; use plain REINFORCE with group advantages
& Our finding provides the mechanistic explanation for their empirical result: clipping is removable because it was already absent.
  Removing a no-op preserves performance by construction. \\
\addlinespace
\citet{chen2025iclr}
& Analyze clipping as entropy regularizer; recommend removal
& Their conclusion is correct in the fresh-log-prob regime (Cause~1 present).
  However, under stale log-probs where clipping becomes active, we find it is essential: our no-clip ablation with stale log-probs collapses within 50 rollouts (Section~\ref{sec:factorial}). \\
\bottomrule
\end{tabular}
\end{table*}

The convergent pattern across these works is notable.
Each paper identifies a problem with GRPO's clipping---it does not prevent entropy collapse (DAPO), it operates at the wrong granularity (GSPO), it creates dead zones (CFPO), it misallocates gradients (ASPO), or it is simply unnecessary (RGRA)---and proposes an independent fix.
Our diagnosis reveals these as facets of a single phenomenon: the importance ratio is architecturally constrained to near unity, so the clip operator is never binding.
Under this lens, the modifications fall into two categories: those that \emph{work around} inert clipping by introducing alternative regularization (GSPO, CFPO), and those that \emph{acknowledge} its inertness by simplifying the objective (RGRA, \citealt{chen2025iclr}).
Neither category addresses the upstream cause---the disconnection between the importance ratio and data staleness.

This suggests that a more principled path forward is to \emph{reconnect} the importance ratio to the data-generating policy, making clipping functional as originally intended.
We pursue this direction in Section~\ref{sec:activating-clipping}, where we show that the combination of stale rollout log-probabilities and tight clipping produces a qualitatively different training regime---one in which clipping actively promotes exploration rather than serving as a dormant safety constraint.


\section{The Inert Clipping Problem}
\label{sec:inert-clipping}

PPO's clipped surrogate objective is the central trust-region mechanism inherited by GRPO.
In this section, we demonstrate that this mechanism is empirically inert across all standard training configurations we tested, identify two independent root causes, and show how this diagnosis unifies a series of recent algorithmic redesigns.

\subsection{Measuring Clipping Activity}
\label{sec:measuring-clipping}

We instrument the PPO loss to track \texttt{pg\_clipfrac}, the fraction of tokens at which the clipped and unclipped objectives differ---i.e., the fraction of tokens where clipping actively modifies the gradient.
A value of $\texttt{pg\_clipfrac} = 0$ means the clip operator in Equation~\ref{eq:ppo-clip} is never binding: for every token, $\text{clip}(r_t, 1-\varepsilon, 1+\varepsilon) \cdot \hat{A}_t = r_t \cdot \hat{A}_t$, and the trust-region constraint contributes nothing to the optimization.

\paragraph{Setup.}
We train Qwen2.5-Math-1.5B on the DAPO-Math-17K dataset under standard GRPO configurations, resuming from a checkpoint at iteration~799.
All runs use the default clip threshold $\varepsilon = 0.2$.
We additionally track the Effective Sample Size (ESS), defined as
\begin{equation}
\label{eq:ess}
\text{ESS} = \frac{1}{\sum_{i} w_i^2}, \quad w_i = \frac{r_i}{\sum_j r_j},
\end{equation}
where $r_i = \pi_\theta(a_i \mid s) / \pi_{\text{old}}(a_i \mid s)$ are the per-token importance ratios.
ESS measures how many effective independent samples remain after importance weighting; $\text{ESS} \approx 1.0$ indicates near-perfect on-policy training.

\paragraph{PPO epochs do not activate clipping.}
At $\text{lr} = 10^{-6}$, we vary the number of PPO epochs $K \in \{1, 3, 5\}$.
In standard GRPO, multiple PPO epochs recompute log-probabilities at the start of each epoch, resetting the importance ratio denominator.
Across all three settings, $\texttt{pg\_clipfrac} = 0.0$ exactly, and $\text{ESS} = 0.999$.
Increasing the number of optimization passes over the same rollout data does not induce any clipping activity.

\paragraph{Data reuse does not activate clipping.}
We increase the learning rate to $10^{-5}$ and set \texttt{--num-steps-per-rollout}$\,{=}\,8$, which partitions each rollout into 8 non-overlapping training chunks processed sequentially \emph{without} recomputing log-probabilities between steps.
This is the most aggressive data-reuse configuration in our experimental grid.
Even here, $\texttt{pg\_clipfrac} = 0.0$ at every step position, and ESS remains above $0.999$.

\paragraph{Comparison with prior observations.}
Our finding of exactly zero clipping is consistent with, but more extreme than, prior reports.
\citet{chen2025iclr} measure $\texttt{pg\_clipfrac} < 0.2\%$ on Qwen2.5-Math-7B at $\text{lr} = 5 \times 10^{-7}$, attributing the low rate to conservative hyperparameters.
\citet{prosperity2025} report $\texttt{pg\_clipfrac} = 0.05\%$ at staleness $s = 0$ in their multi-policy framework.
Neither work investigates why the rate is near zero or identifies the architectural cause.
The small residual clipping they observe likely reflects numerical differences in their pipelines (larger models, different checkpoint states, or floating-point accumulation effects), but the qualitative conclusion is the same: \emph{clipping is essentially inactive in standard GRPO training}.


\subsection{Root Cause: Two Independent Mechanisms}
\label{sec:root-cause}

We identify two independent mechanisms that jointly render PPO clipping inert.
Their independence is important: removing either one alone is sufficient to produce nonzero clipping activity.

\subsubsection{Cause 1: Architectural Disconnection}
\label{sec:cause-architectural}

The PPO clipped surrogate objective computes the importance ratio as
\begin{equation}
\label{eq:importance-ratio}
r_t(\theta) = \frac{\pi_\theta(a_t \mid s_t)}{\pi_{\text{old}}(a_t \mid s_t)},
\end{equation}
where, in classical PPO, $\pi_{\text{old}}$ represents the \emph{behavior policy} that generated the data.
In GRPO as implemented in standard training frameworks, the denominator uses \texttt{batch["log\_probs"]}---log-probabilities \emph{recomputed} from the current model at the start of each training step---rather than \texttt{batch["rollout\_log\_probs"]}, the log-probabilities recorded at rollout generation time.

This architectural choice means that both numerator and denominator of $r_t(\theta)$ track the same policy (or very nearly so, differing only by the gradient updates within a single step).
Consequently, $r_t(\theta) \approx 1$ regardless of how far the current policy $\pi_\theta$ has drifted from the policy $\pi_{\text{rollout}}$ that actually generated the training data.
The importance ratio is \emph{disconnected} from data staleness: the PPO loss cannot detect, and therefore cannot constrain, policy drift relative to the data-generating distribution.

\subsubsection{Cause 2: Conservative Learning Rates}
\label{sec:cause-lr}

Standard GRPO training uses learning rates in the range $10^{-6}$ to $10^{-5}$.
At these rates, the per-step change in log-probabilities is small enough that even a correctly computed importance ratio (using stale rollout log-probs) would rarely exceed the clip threshold $1 + \varepsilon = 1.2$.
This provides a second, independent barrier to clipping activation: the policy simply does not move fast enough within a single rollout's training steps for the trust region to bind.

\subsubsection{Evidence for Independence}
\label{sec:evidence-independence}

We isolate the two causes by introducing a \texttt{--use-rollout-logprobs} flag that replaces the denominator of $r_t(\theta)$ with the stale rollout log-probabilities, removing Cause~1 while leaving Cause~2 intact.
At $\text{lr} = 10^{-5}$ with \texttt{--num-steps-per-rollout}$\,{=}\,8$:

\begin{itemize}
    \item \textbf{With stale log-probs} (Cause~1 removed): We observe a clear ESS gradient within each rollout, declining from $0.999$ at step~1 to $0.988$ at step~8.
    The variable \texttt{pg\_clipfrac} is nonzero at \emph{all} step positions, rising from $0.001$ (step~1) to $0.024$ (step~8) as accumulated policy drift pushes more tokens past the clip threshold.
    Over 218 rollouts, \texttt{pg\_clipfrac} is consistently positive.

    \item \textbf{Without stale log-probs} (Cause~1 present): Under identical hyperparameters, $\texttt{pg\_clipfrac} = 0.0$ at every step position regardless of $N$.
    The recomputed denominator masks all policy drift.
\end{itemize}

This contrast establishes that Cause~1 and Cause~2 are independent.
Removing the architectural disconnection (Cause~1) is \emph{sufficient} to activate clipping even at the conservative learning rate of $10^{-5}$---the rate merely modulates the \emph{magnitude} of clipping activity, not its presence.
Conversely, increasing the learning rate alone (without fixing the log-prob reference) cannot activate clipping, because the recomputed denominator absorbs all policy changes.


\subsection{Implications for Recent Algorithmic Redesigns}
\label{sec:implications}

Our diagnosis provides a unifying explanation for a striking pattern in recent work: at least six independent papers have redesigned or removed GRPO's clipping mechanism, each motivated by different empirical observations, yet all responding to the same underlying pathology.
Table~\ref{tab:related-redesigns} summarizes how each algorithm's design choice relates to our two-cause diagnosis.

\begin{table*}[t]
\centering
\caption{Recent clipping redesigns interpreted through the lens of inert clipping.
Each modification addresses a symptom of the same root cause: the PPO importance ratio is architecturally near unity, rendering clipping inactive.
``Fresh log-probs'' indicates the algorithm retains the recomputed denominator (Cause~1 unresolved).}
\label{tab:related-redesigns}
\small
\begin{tabular}{@{}p{2.0cm}p{4.5cm}p{7.5cm}@{}}
\toprule
\textbf{Algorithm} & \textbf{Modification} & \textbf{Our Explanation} \\
\midrule
DAPO \newline \citep{dapo2025}
& Asymmetric clip bounds ($\varepsilon_{\text{low}} = 0.2$, $\varepsilon_{\text{high}} = 0.28$)
& Presupposes active clipping.
  If $\texttt{pg\_clipfrac} = 0$, Clip-Higher has no tokens to asymmetrically clip.
  Their reported gains (+8 AIME) are attributable to Dynamic Sampling, not the clip modification (+2 AIME). \\
\addlinespace
GSPO \newline \citep{gspo2025}
& Sequence-level importance sampling replacing token-level
& Retains fresh log-probs; Cause~1 persists.
  Their finding that GSPO clips $100\times$ more tokens than standard GRPO directly supports our diagnosis: GRPO clips almost nothing. \\
\addlinespace
CFPO \newline \citep{cfpo2025}
& Quadratic TV-divergence penalty replacing clip
& Their penalty term is also approximately zero when $r_t \approx 1$.
  The zero-gradient problem they address only manifests after ratios deviate from unity---which requires resolving our upstream Cause~1 first. \\
\addlinespace
ASPO \newline \citep{aspo2025}
& Inverted IS ratios for positive-advantage tokens
& They report average IS weight $\approx 1.0004$ but do not investigate this observation.
  Their gradient-direction fix has negligible magnitude when $r_t \approx 1$; the effective update is dominated by the advantage term alone. \\
\addlinespace
RGRA \newline \citep{rgra2025}
& Remove clipping entirely; use plain REINFORCE with group advantages
& Our finding provides the mechanistic explanation for their empirical result: clipping is removable because it was already absent.
  Removing a no-op preserves performance by construction. \\
\addlinespace
\citet{chen2025iclr}
& Analyze clipping as entropy regularizer; recommend removal
& Their conclusion is correct in the fresh-log-prob regime (Cause~1 present).
  However, under stale log-probs where clipping becomes active, we find it is essential: our no-clip ablation with stale log-probs collapses within 50 rollouts (Section~\ref{sec:factorial}). \\
\bottomrule
\end{tabular}
\end{table*}

The convergent pattern across these works is notable.
Each paper identifies a problem with GRPO's clipping---it does not prevent entropy collapse (DAPO), it operates at the wrong granularity (GSPO), it creates dead zones (CFPO), it misallocates gradients (ASPO), or it is simply unnecessary (RGRA)---and proposes an independent fix.
Our diagnosis reveals these as facets of a single phenomenon: the importance ratio is architecturally constrained to near unity, so the clip operator is never binding.
Under this lens, the modifications fall into two categories: those that \emph{work around} inert clipping by introducing alternative regularization (GSPO, CFPO), and those that \emph{acknowledge} its inertness by simplifying the objective (RGRA, \citealt{chen2025iclr}).
Neither category addresses the upstream cause---the disconnection between the importance ratio and data staleness.

This suggests that a more principled path forward is to \emph{reconnect} the importance ratio to the data-generating policy, making clipping functional as originally intended.
We pursue this direction in Section~\ref{sec:activating-clipping}, where we show that the combination of stale rollout log-probabilities and tight clipping produces a qualitatively different training regime---one in which clipping actively promotes exploration rather than serving as a dormant safety constraint.


\section{The Inert Clipping Problem}
\label{sec:inert-clipping}

PPO's clipped surrogate objective is the central trust-region mechanism inherited by GRPO.
In this section, we demonstrate that this mechanism is empirically inert across all standard training configurations we tested, identify two independent root causes, and show how this diagnosis unifies a series of recent algorithmic redesigns.

\subsection{Measuring Clipping Activity}
\label{sec:measuring-clipping}

We instrument the PPO loss to track \texttt{pg\_clipfrac}, the fraction of tokens at which the clipped and unclipped objectives differ---i.e., the fraction of tokens where clipping actively modifies the gradient.
A value of $\texttt{pg\_clipfrac} = 0$ means the clip operator in Equation~\ref{eq:ppo-clip} is never binding: for every token, $\text{clip}(r_t, 1-\varepsilon, 1+\varepsilon) \cdot \hat{A}_t = r_t \cdot \hat{A}_t$, and the trust-region constraint contributes nothing to the optimization.

\paragraph{Setup.}
We train Qwen2.5-Math-1.5B on the DAPO-Math-17K dataset under standard GRPO configurations, resuming from a checkpoint at iteration~799.
All runs use the default clip threshold $\varepsilon = 0.2$.
We additionally track the Effective Sample Size (ESS), defined as
\begin{equation}
\label{eq:ess}
\text{ESS} = \frac{1}{\sum_{i} w_i^2}, \quad w_i = \frac{r_i}{\sum_j r_j},
\end{equation}
where $r_i = \pi_\theta(a_i \mid s) / \pi_{\text{old}}(a_i \mid s)$ are the per-token importance ratios.
ESS measures how many effective independent samples remain after importance weighting; $\text{ESS} \approx 1.0$ indicates near-perfect on-policy training.

\paragraph{PPO epochs do not activate clipping.}
At $\text{lr} = 10^{-6}$, we vary the number of PPO epochs $K \in \{1, 3, 5\}$.
In standard GRPO, multiple PPO epochs recompute log-probabilities at the start of each epoch, resetting the importance ratio denominator.
Across all three settings, $\texttt{pg\_clipfrac} = 0.0$ exactly, and $\text{ESS} = 0.999$.
Increasing the number of optimization passes over the same rollout data does not induce any clipping activity.

\paragraph{Data reuse does not activate clipping.}
We increase the learning rate to $10^{-5}$ and set \texttt{--num-steps-per-rollout}$\,{=}\,8$, which partitions each rollout into 8 non-overlapping training chunks processed sequentially \emph{without} recomputing log-probabilities between steps.
This is the most aggressive data-reuse configuration in our experimental grid.
Even here, $\texttt{pg\_clipfrac} = 0.0$ at every step position, and ESS remains above $0.999$.

\paragraph{Comparison with prior observations.}
Our finding of exactly zero clipping is consistent with, but more extreme than, prior reports.
\citet{chen2025iclr} measure $\texttt{pg\_clipfrac} < 0.2\%$ on Qwen2.5-Math-7B at $\text{lr} = 5 \times 10^{-7}$, attributing the low rate to conservative hyperparameters.
\citet{prosperity2025} report $\texttt{pg\_clipfrac} = 0.05\%$ at staleness $s = 0$ in their multi-policy framework.
Neither work investigates why the rate is near zero or identifies the architectural cause.
The small residual clipping they observe likely reflects numerical differences in their pipelines (larger models, different checkpoint states, or floating-point accumulation effects), but the qualitative conclusion is the same: \emph{clipping is essentially inactive in standard GRPO training}.


\subsection{Root Cause: Two Independent Mechanisms}
\label{sec:root-cause}

We identify two independent mechanisms that jointly render PPO clipping inert.
Their independence is important: removing either one alone is sufficient to produce nonzero clipping activity.

\subsubsection{Cause 1: Architectural Disconnection}
\label{sec:cause-architectural}

The PPO clipped surrogate objective computes the importance ratio as
\begin{equation}
\label{eq:importance-ratio}
r_t(\theta) = \frac{\pi_\theta(a_t \mid s_t)}{\pi_{\text{old}}(a_t \mid s_t)},
\end{equation}
where, in classical PPO, $\pi_{\text{old}}$ represents the \emph{behavior policy} that generated the data.
In GRPO as implemented in standard training frameworks, the denominator uses \texttt{batch["log\_probs"]}---log-probabilities \emph{recomputed} from the current model at the start of each training step---rather than \texttt{batch["rollout\_log\_probs"]}, the log-probabilities recorded at rollout generation time.

This architectural choice means that both numerator and denominator of $r_t(\theta)$ track the same policy (or very nearly so, differing only by the gradient updates within a single step).
Consequently, $r_t(\theta) \approx 1$ regardless of how far the current policy $\pi_\theta$ has drifted from the policy $\pi_{\text{rollout}}$ that actually generated the training data.
The importance ratio is \emph{disconnected} from data staleness: the PPO loss cannot detect, and therefore cannot constrain, policy drift relative to the data-generating distribution.

\subsubsection{Cause 2: Conservative Learning Rates}
\label{sec:cause-lr}

Standard GRPO training uses learning rates in the range $10^{-6}$ to $10^{-5}$.
At these rates, the per-step change in log-probabilities is small enough that even a correctly computed importance ratio (using stale rollout log-probs) would rarely exceed the clip threshold $1 + \varepsilon = 1.2$.
This provides a second, independent barrier to clipping activation: the policy simply does not move fast enough within a single rollout's training steps for the trust region to bind.

\subsubsection{Evidence for Independence}
\label{sec:evidence-independence}

We isolate the two causes by introducing a \texttt{--use-rollout-logprobs} flag that replaces the denominator of $r_t(\theta)$ with the stale rollout log-probabilities, removing Cause~1 while leaving Cause~2 intact.
At $\text{lr} = 10^{-5}$ with \texttt{--num-steps-per-rollout}$\,{=}\,8$:

\begin{itemize}
    \item \textbf{With stale log-probs} (Cause~1 removed): We observe a clear ESS gradient within each rollout, declining from $0.999$ at step~1 to $0.988$ at step~8.
    The variable \texttt{pg\_clipfrac} is nonzero at \emph{all} step positions, rising from $0.001$ (step~1) to $0.024$ (step~8) as accumulated policy drift pushes more tokens past the clip threshold.
    Over 218 rollouts, \texttt{pg\_clipfrac} is consistently positive.

    \item \textbf{Without stale log-probs} (Cause~1 present): Under identical hyperparameters, $\texttt{pg\_clipfrac} = 0.0$ at every step position regardless of $N$.
    The recomputed denominator masks all policy drift.
\end{itemize}

This contrast establishes that Cause~1 and Cause~2 are independent.
Removing the architectural disconnection (Cause~1) is \emph{sufficient} to activate clipping even at the conservative learning rate of $10^{-5}$---the rate merely modulates the \emph{magnitude} of clipping activity, not its presence.
Conversely, increasing the learning rate alone (without fixing the log-prob reference) cannot activate clipping, because the recomputed denominator absorbs all policy changes.


\subsection{Implications for Recent Algorithmic Redesigns}
\label{sec:implications}

Our diagnosis provides a unifying explanation for a striking pattern in recent work: at least six independent papers have redesigned or removed GRPO's clipping mechanism, each motivated by different empirical observations, yet all responding to the same underlying pathology.
Table~\ref{tab:related-redesigns} summarizes how each algorithm's design choice relates to our two-cause diagnosis.

\begin{table*}[t]
\centering
\caption{Recent clipping redesigns interpreted through the lens of inert clipping.
Each modification addresses a symptom of the same root cause: the PPO importance ratio is architecturally near unity, rendering clipping inactive.
``Fresh log-probs'' indicates the algorithm retains the recomputed denominator (Cause~1 unresolved).}
\label{tab:related-redesigns}
\small
\begin{tabular}{@{}p{2.0cm}p{4.5cm}p{7.5cm}@{}}
\toprule
\textbf{Algorithm} & \textbf{Modification} & \textbf{Our Explanation} \\
\midrule
DAPO \newline \citep{dapo2025}
& Asymmetric clip bounds ($\varepsilon_{\text{low}} = 0.2$, $\varepsilon_{\text{high}} = 0.28$)
& Presupposes active clipping.
  If $\texttt{pg\_clipfrac} = 0$, Clip-Higher has no tokens to asymmetrically clip.
  Their reported gains (+8 AIME) are attributable to Dynamic Sampling, not the clip modification (+2 AIME). \\
\addlinespace
GSPO \newline \citep{gspo2025}
& Sequence-level importance sampling replacing token-level
& Retains fresh log-probs; Cause~1 persists.
  Their finding that GSPO clips $100\times$ more tokens than standard GRPO directly supports our diagnosis: GRPO clips almost nothing. \\
\addlinespace
CFPO \newline \citep{cfpo2025}
& Quadratic TV-divergence penalty replacing clip
& Their penalty term is also approximately zero when $r_t \approx 1$.
  The zero-gradient problem they address only manifests after ratios deviate from unity---which requires resolving our upstream Cause~1 first. \\
\addlinespace
ASPO \newline \citep{aspo2025}
& Inverted IS ratios for positive-advantage tokens
& They report average IS weight $\approx 1.0004$ but do not investigate this observation.
  Their gradient-direction fix has negligible magnitude when $r_t \approx 1$; the effective update is dominated by the advantage term alone. \\
\addlinespace
RGRA \newline \citep{rgra2025}
& Remove clipping entirely; use plain REINFORCE with group advantages
& Our finding provides the mechanistic explanation for their empirical result: clipping is removable because it was already absent.
  Removing a no-op preserves performance by construction. \\
\addlinespace
\citet{chen2025iclr}
& Analyze clipping as entropy regularizer; recommend removal
& Their conclusion is correct in the fresh-log-prob regime (Cause~1 present).
  However, under stale log-probs where clipping becomes active, we find it is essential: our no-clip ablation with stale log-probs collapses within 50 rollouts (Section~\ref{sec:factorial}). \\
\bottomrule
\end{tabular}
\end{table*}

The convergent pattern across these works is notable.
Each paper identifies a problem with GRPO's clipping---it does not prevent entropy collapse (DAPO), it operates at the wrong granularity (GSPO), it creates dead zones (CFPO), it misallocates gradients (ASPO), or it is simply unnecessary (RGRA)---and proposes an independent fix.
Our diagnosis reveals these as facets of a single phenomenon: the importance ratio is architecturally constrained to near unity, so the clip operator is never binding.
Under this lens, the modifications fall into two categories: those that \emph{work around} inert clipping by introducing alternative regularization (GSPO, CFPO), and those that \emph{acknowledge} its inertness by simplifying the objective (RGRA, \citealt{chen2025iclr}).
Neither category addresses the upstream cause---the disconnection between the importance ratio and data staleness.

This suggests that a more principled path forward is to \emph{reconnect} the importance ratio to the data-generating policy, making clipping functional as originally intended.
We pursue this direction in Section~\ref{sec:activating-clipping}, where we show that the combination of stale rollout log-probabilities and tight clipping produces a qualitatively different training regime---one in which clipping actively promotes exploration rather than serving as a dormant safety constraint.


\section{The Inert Clipping Problem}
\label{sec:inert-clipping}

PPO's clipped surrogate objective is the central trust-region mechanism inherited by GRPO.
In this section, we demonstrate that this mechanism is empirically inert across all standard training configurations we tested, identify two independent root causes, and show how this diagnosis unifies a series of recent algorithmic redesigns.

\subsection{Measuring Clipping Activity}
\label{sec:measuring-clipping}

We instrument the PPO loss to track \texttt{pg\_clipfrac}, the fraction of tokens at which the clipped and unclipped objectives differ---i.e., the fraction of tokens where clipping actively modifies the gradient.
A value of $\texttt{pg\_clipfrac} = 0$ means the clip operator in Equation~\ref{eq:ppo-clip} is never binding: for every token, $\text{clip}(r_t, 1-\varepsilon, 1+\varepsilon) \cdot \hat{A}_t = r_t \cdot \hat{A}_t$, and the trust-region constraint contributes nothing to the optimization.

\paragraph{Setup.}
We train Qwen2.5-Math-1.5B on the DAPO-Math-17K dataset under standard GRPO configurations, resuming from a checkpoint at iteration~799.
All runs use the default clip threshold $\varepsilon = 0.2$.
We additionally track the Effective Sample Size (ESS), defined as
\begin{equation}
\label{eq:ess}
\text{ESS} = \frac{1}{\sum_{i} w_i^2}, \quad w_i = \frac{r_i}{\sum_j r_j},
\end{equation}
where $r_i = \pi_\theta(a_i \mid s) / \pi_{\text{old}}(a_i \mid s)$ are the per-token importance ratios.
ESS measures how many effective independent samples remain after importance weighting; $\text{ESS} \approx 1.0$ indicates near-perfect on-policy training.

\paragraph{PPO epochs do not activate clipping.}
At $\text{lr} = 10^{-6}$, we vary the number of PPO epochs $K \in \{1, 3, 5\}$.
In standard GRPO, multiple PPO epochs recompute log-probabilities at the start of each epoch, resetting the importance ratio denominator.
Across all three settings, $\texttt{pg\_clipfrac} = 0.0$ exactly, and $\text{ESS} = 0.999$.
Increasing the number of optimization passes over the same rollout data does not induce any clipping activity.

\paragraph{Data reuse does not activate clipping.}
We increase the learning rate to $10^{-5}$ and set \texttt{--num-steps-per-rollout}$\,{=}\,8$, which partitions each rollout into 8 non-overlapping training chunks processed sequentially \emph{without} recomputing log-probabilities between steps.
This is the most aggressive data-reuse configuration in our experimental grid.
Even here, $\texttt{pg\_clipfrac} = 0.0$ at every step position, and ESS remains above $0.999$.

\paragraph{Comparison with prior observations.}
Our finding of exactly zero clipping is consistent with, but more extreme than, prior reports.
\citet{chen2025iclr} measure $\texttt{pg\_clipfrac} < 0.2\%$ on Qwen2.5-Math-7B at $\text{lr} = 5 \times 10^{-7}$, attributing the low rate to conservative hyperparameters.
\citet{prosperity2025} report $\texttt{pg\_clipfrac} = 0.05\%$ at staleness $s = 0$ in their multi-policy framework.
Neither work investigates why the rate is near zero or identifies the architectural cause.
The small residual clipping they observe likely reflects numerical differences in their pipelines (larger models, different checkpoint states, or floating-point accumulation effects), but the qualitative conclusion is the same: \emph{clipping is essentially inactive in standard GRPO training}.


\subsection{Root Cause: Two Independent Mechanisms}
\label{sec:root-cause}

We identify two independent mechanisms that jointly render PPO clipping inert.
Their independence is important: removing either one alone is sufficient to produce nonzero clipping activity.

\subsubsection{Cause 1: Architectural Disconnection}
\label{sec:cause-architectural}

The PPO clipped surrogate objective computes the importance ratio as
\begin{equation}
\label{eq:importance-ratio}
r_t(\theta) = \frac{\pi_\theta(a_t \mid s_t)}{\pi_{\text{old}}(a_t \mid s_t)},
\end{equation}
where, in classical PPO, $\pi_{\text{old}}$ represents the \emph{behavior policy} that generated the data.
In GRPO as implemented in standard training frameworks, the denominator uses \texttt{batch["log\_probs"]}---log-probabilities \emph{recomputed} from the current model at the start of each training step---rather than \texttt{batch["rollout\_log\_probs"]}, the log-probabilities recorded at rollout generation time.

This architectural choice means that both numerator and denominator of $r_t(\theta)$ track the same policy (or very nearly so, differing only by the gradient updates within a single step).
Consequently, $r_t(\theta) \approx 1$ regardless of how far the current policy $\pi_\theta$ has drifted from the policy $\pi_{\text{rollout}}$ that actually generated the training data.
The importance ratio is \emph{disconnected} from data staleness: the PPO loss cannot detect, and therefore cannot constrain, policy drift relative to the data-generating distribution.

\subsubsection{Cause 2: Conservative Learning Rates}
\label{sec:cause-lr}

Standard GRPO training uses learning rates in the range $10^{-6}$ to $10^{-5}$.
At these rates, the per-step change in log-probabilities is small enough that even a correctly computed importance ratio (using stale rollout log-probs) would rarely exceed the clip threshold $1 + \varepsilon = 1.2$.
This provides a second, independent barrier to clipping activation: the policy simply does not move fast enough within a single rollout's training steps for the trust region to bind.

\subsubsection{Evidence for Independence}
\label{sec:evidence-independence}

We isolate the two causes by introducing a \texttt{--use-rollout-logprobs} flag that replaces the denominator of $r_t(\theta)$ with the stale rollout log-probabilities, removing Cause~1 while leaving Cause~2 intact.
At $\text{lr} = 10^{-5}$ with \texttt{--num-steps-per-rollout}$\,{=}\,8$:

\begin{itemize}
    \item \textbf{With stale log-probs} (Cause~1 removed): We observe a clear ESS gradient within each rollout, declining from $0.999$ at step~1 to $0.988$ at step~8.
    The variable \texttt{pg\_clipfrac} is nonzero at \emph{all} step positions, rising from $0.001$ (step~1) to $0.024$ (step~8) as accumulated policy drift pushes more tokens past the clip threshold.
    Over 218 rollouts, \texttt{pg\_clipfrac} is consistently positive.

    \item \textbf{Without stale log-probs} (Cause~1 present): Under identical hyperparameters, $\texttt{pg\_clipfrac} = 0.0$ at every step position regardless of $N$.
    The recomputed denominator masks all policy drift.
\end{itemize}

This contrast establishes that Cause~1 and Cause~2 are independent.
Removing the architectural disconnection (Cause~1) is \emph{sufficient} to activate clipping even at the conservative learning rate of $10^{-5}$---the rate merely modulates the \emph{magnitude} of clipping activity, not its presence.
Conversely, increasing the learning rate alone (without fixing the log-prob reference) cannot activate clipping, because the recomputed denominator absorbs all policy changes.


\subsection{Implications for Recent Algorithmic Redesigns}
\label{sec:implications}

Our diagnosis provides a unifying explanation for a striking pattern in recent work: at least six independent papers have redesigned or removed GRPO's clipping mechanism, each motivated by different empirical observations, yet all responding to the same underlying pathology.
Table~\ref{tab:related-redesigns} summarizes how each algorithm's design choice relates to our two-cause diagnosis.

\begin{table*}[t]
\centering
\caption{Recent clipping redesigns interpreted through the lens of inert clipping.
Each modification addresses a symptom of the same root cause: the PPO importance ratio is architecturally near unity, rendering clipping inactive.
``Fresh log-probs'' indicates the algorithm retains the recomputed denominator (Cause~1 unresolved).}
\label{tab:related-redesigns}
\small
\begin{tabular}{@{}p{2.0cm}p{4.5cm}p{7.5cm}@{}}
\toprule
\textbf{Algorithm} & \textbf{Modification} & \textbf{Our Explanation} \\
\midrule
DAPO \newline \citep{dapo2025}
& Asymmetric clip bounds ($\varepsilon_{\text{low}} = 0.2$, $\varepsilon_{\text{high}} = 0.28$)
& Presupposes active clipping.
  If $\texttt{pg\_clipfrac} = 0$, Clip-Higher has no tokens to asymmetrically clip.
  Their reported gains (+8 AIME) are attributable to Dynamic Sampling, not the clip modification (+2 AIME). \\
\addlinespace
GSPO \newline \citep{gspo2025}
& Sequence-level importance sampling replacing token-level
& Retains fresh log-probs; Cause~1 persists.
  Their finding that GSPO clips $100\times$ more tokens than standard GRPO directly supports our diagnosis: GRPO clips almost nothing. \\
\addlinespace
CFPO \newline \citep{cfpo2025}
& Quadratic TV-divergence penalty replacing clip
& Their penalty term is also approximately zero when $r_t \approx 1$.
  The zero-gradient problem they address only manifests after ratios deviate from unity---which requires resolving our upstream Cause~1 first. \\
\addlinespace
ASPO \newline \citep{aspo2025}
& Inverted IS ratios for positive-advantage tokens
& They report average IS weight $\approx 1.0004$ but do not investigate this observation.
  Their gradient-direction fix has negligible magnitude when $r_t \approx 1$; the effective update is dominated by the advantage term alone. \\
\addlinespace
RGRA \newline \citep{rgra2025}
& Remove clipping entirely; use plain REINFORCE with group advantages
& Our finding provides the mechanistic explanation for their empirical result: clipping is removable because it was already absent.
  Removing a no-op preserves performance by construction. \\
\addlinespace
\citet{chen2025iclr}
& Analyze clipping as entropy regularizer; recommend removal
& Their conclusion is correct in the fresh-log-prob regime (Cause~1 present).
  However, under stale log-probs where clipping becomes active, we find it is essential: our no-clip ablation with stale log-probs collapses within 50 rollouts (Section~\ref{sec:factorial}). \\
\bottomrule
\end{tabular}
\end{table*}

The convergent pattern across these works is notable.
Each paper identifies a problem with GRPO's clipping---it does not prevent entropy collapse (DAPO), it operates at the wrong granularity (GSPO), it creates dead zones (CFPO), it misallocates gradients (ASPO), or it is simply unnecessary (RGRA)---and proposes an independent fix.
Our diagnosis reveals these as facets of a single phenomenon: the importance ratio is architecturally constrained to near unity, so the clip operator is never binding.
Under this lens, the modifications fall into two categories: those that \emph{work around} inert clipping by introducing alternative regularization (GSPO, CFPO), and those that \emph{acknowledge} its inertness by simplifying the objective (RGRA, \citealt{chen2025iclr}).
Neither category addresses the upstream cause---the disconnection between the importance ratio and data staleness.

This suggests that a more principled path forward is to \emph{reconnect} the importance ratio to the data-generating policy, making clipping functional as originally intended.
We pursue this direction in Section~\ref{sec:activating-clipping}, where we show that the combination of stale rollout log-probabilities and tight clipping produces a qualitatively different training regime---one in which clipping actively promotes exploration rather than serving as a dormant safety constraint.
